% Aula 02 --- o melhor preditor linear (o "oráculo") quando r não é linear.
Circula por aí a ideia de que a regressão linear "só vale se o modelo estiver
correto". Vimos durante o curso que isso é impreciso: o MQO converge para uma função
bem definida mesmo quando $r$ não é linear.

Essa função é o \textbf{melhor preditor linear},
\[
  \bbeta^\ast = \argmin_{\bbeta} \E\big[(Y-\bbeta^\top\X)^2\big],
\]
que existe sempre que $\Sigma=\E[\X\X^\top]$ é bem definida e invertível ---
\emph{sem} supor nada sobre a forma de $r$.
\begin{itens}
  \item \pts{0,9} Derive $\bbeta^\ast=\Sigma^{-1}\bm\alpha$, onde
        $\bm\alpha=\E[Y\X]$. É o mesmo cálculo das equações normais, com
        esperanças no lugar de somas.
  \item \pts{1,0} O estimador de mínimos quadrados pode ser escrito como
        $\widehat\bbeta=\widehat\Sigma^{-1}\widehat{\bm\alpha}$, com
        $\widehat\Sigma=\frac1n\sum_i\X_i\X_i^\top$ e
        $\widehat{\bm\alpha}=\frac1n\sum_i Y_i\X_i$. Olhe para essas duas
        expressões e diga o que elas são, estatisticamente. Em seguida, explique
        por que $\widehat\bbeta\to\bbeta^\ast$ em probabilidade sem que se precise
        supor linearidade em lugar nenhum.
  \item \pts{1,0} Todo teorema garante uma coisa e deixa de garantir outra.
        Enuncie as duas, com cuidado. Em particular: o resultado diz que o risco
        de $\widehat\bbeta$ é pequeno?
  \item \pts{1,1} Um colega ajusta uma regressão linear a dados em que $r$ é
        visivelmente curva e conclui: "o coeficiente deu $2{,}3$, logo aumentar
        $x_1$ em uma unidade aumenta $Y$ em $2{,}3$". Separe o que é leitura
        legítima de $\bbeta^\ast$ do que não é.
\end{itens}

\begin{solucao}
Hipóteses em uso: $\E[Y^2]<\infty$ e $\E\norm{\X}^2<\infty$, para que $\Sigma$ e
$\bm\alpha$ existam (a existência de $\bm\alpha$ sai por Cauchy--Schwarz), e $\Sigma$
invertível. Como $\bm v^\top\Sigma\bm v=\E\big[(\bm v^\top\X)^2\big]\ge0$ para todo
$\bm v$, $\Sigma$ é semidefinida positiva; sendo invertível, é definida positiva.

\smallskip
\textbf{(a)} Expanda o objetivo; as esperanças são números fixos:
\[
  Q(\bbeta)=\E\big[(Y-\bbeta^\top\X)^2\big]
  =\E[Y^2]-2\,\bbeta^\top\underbrace{\E[Y\X]}_{\bm\alpha}
   +\bbeta^\top\underbrace{\E[\X\X^\top]}_{\Sigma}\bbeta .
\]
É uma função quadrática de $\bbeta$, com gradiente $\nabla Q(\bbeta)=-2\bm\alpha+2\Sigma\bbeta$
e Hessiana $2\Sigma$, definida positiva. Logo $Q$ é estritamente convexa, e seu único
ponto crítico é o mínimo global:
\[
  \Sigma\bbeta=\bm\alpha
  \iff
  \bbeta^\ast=\Sigma^{-1}\bm\alpha .
\]
São as equações normais da aula com $\E[\X\X^\top]$ no lugar de
$\frac1n\mathbb{X}^\top\mathbb{X}$ e $\E[Y\X]$ no lugar de $\frac1n\mathbb{X}^\top\bm y$:
o mesmo problema, escrito na população em vez de na amostra.

Vale conhecer uma segunda demonstração, sem derivadas, porque ela rende uma identidade
que usaremos em (c). Some e subtraia $\bbeta^{\ast\top}\X$ dentro do quadrado:
\[
\begin{aligned}
  Q(\bbeta)
  &=\E\Big[\big((Y-\bbeta^{\ast\top}\X)-(\bbeta-\bbeta^\ast)^\top\X\big)^2\Big]\\
  &=Q(\bbeta^\ast)
   -2(\bbeta-\bbeta^\ast)^\top\underbrace{\E\big[\X\,(Y-\X^\top\bbeta^\ast)\big]}_{=\;\bm\alpha-\Sigma\bbeta^\ast\;=\;\bm 0}
   +(\bbeta-\bbeta^\ast)^\top\Sigma\,(\bbeta-\bbeta^\ast),
\end{aligned}
\]
ou seja,
\[
  Q(\bbeta)=Q(\bbeta^\ast)+(\bbeta-\bbeta^\ast)^\top\Sigma\,(\bbeta-\bbeta^\ast).
  \tag{$\dagger$}
\]
O último termo é estritamente positivo para $\bbeta\neq\bbeta^\ast$, o que prova de uma
vez que $\bbeta^\ast$ é o mínimo e que é o único. É o mesmo "some e subtraia" do
Exercício~\ref{ex:01-ropt}, e funciona pela mesma razão. Lá, o resíduo $Y-r(\X)$ era
ortogonal a toda função de $\X$; aqui, $\E\big[\X(Y-\X^\top\bbeta^\ast)\big]=\bm0$ diz
que o resíduo do melhor preditor linear é ortogonal a cada covariável. $r(\X)$ é a
projeção de $Y$ sobre todas as funções de $\X$; $\bbeta^{\ast\top}\X$, a projeção sobre
as lineares.

\smallskip
\textbf{(b)} $\widehat\Sigma$ e $\widehat{\bm\alpha}$ são \textbf{médias amostrais} de
variáveis i.i.d.: a primeira, das matrizes $\X_i\X_i^\top$, cuja esperança é $\Sigma$;
a segunda, dos vetores $Y_i\X_i$, cuja esperança é $\bm\alpha$. É essa observação que
resolve o item.

Pela lei dos grandes números, aplicada coordenada a coordenada (as esperanças existem
pelas hipóteses acima), $\widehat\Sigma\to\Sigma$ e $\widehat{\bm\alpha}\to\bm\alpha$ em
probabilidade. A função $(A,\bm a)\mapsto A^{-1}\bm a$ é contínua no conjunto aberto das
matrizes invertíveis, que contém $\Sigma$; por isso $\widehat\Sigma$ é invertível com
probabilidade tendendo a $1$, e o teorema de Mann--Wald (da aplicação contínua)
transporta a convergência:
\[
  \widehat\bbeta=\widehat\Sigma^{-1}\widehat{\bm\alpha}
  \;\xrightarrow{\ P\ }\;\Sigma^{-1}\bm\alpha=\bbeta^\ast .
\]
Em nenhum passo apareceu a forma de $r$. O alvo $\bbeta^\ast$ é definido por um
problema de minimização que faz sentido para qualquer distribuição com segundos
momentos finitos; o que a linearidade de $r$ acrescentaria seria só a coincidência
$\bbeta^{\ast\top}\x=r(\x)$ --- agradável, mas dispensável para a convergência.

\smallskip
\textbf{(c)} Duas identidades do curso dizem tudo. Aplique a do
Exercício~\ref{ex:01-ropt}(a) com $g(\x)=\bbeta^{\ast\top}\x$, e depois $(\dagger)$ com
$\bbeta=\widehat\bbeta$ (fixado o treino, o risco do ajuste é
$\risco(\widehat g)=Q(\widehat\bbeta)$):
\[
  \risco(\widehat g)=
  \underbrace{\E\big[\Var(Y\mid\X)\big]}_{\text{irredutível}}
  +\underbrace{\E\big[(r(\X)-\bbeta^{\ast\top}\X)^2\big]}_{\text{erro de aproximação}}
  +\underbrace{(\widehat\bbeta-\bbeta^\ast)^\top\Sigma\,(\widehat\bbeta-\bbeta^\ast)}_{\text{erro de estimação}} .
\]
Com a decomposição à vista, as duas coisas se enunciam sem ambiguidade.

\textbf{O que o teorema garante:} a terceira parcela vai a zero, porque
$\widehat\bbeta\to\bbeta^\ast$ e a forma quadrática é contínua. Com $n$ grande, o MQO
chega tão perto quanto se queira do melhor que a classe das funções lineares tem a
oferecer, e $\risco(\widehat g)\to\risco(g_{\bbeta^\ast})$.

\textbf{O que ele não garante:} que esse melhor seja bom. A segunda parcela não depende
de $n$, e o teorema não diz nada sobre ela; se $r$ é muito curva, ela é grande, e
convergir para uma aproximação ruim continua sendo ruim. Então não: o resultado
\emph{não} diz que o risco de $\widehat\bbeta$ é pequeno, só que ele se aproxima do
risco do oráculo. Tampouco diz \emph{com que rapidez}: é um resultado de convergência,
sem taxa.

Na linguagem da Aula~01, o teorema cuida do erro de estimação --- a variância, e a
parte do viés que vem de $n$ ser finito --- e deixa intacto o viés de aproximação, que
é o preço fixo de ter escolhido uma classe rígida.

\smallskip
\textbf{(d)} \emph{A leitura legítima.} $2{,}3$ é o coeficiente de $x_1$ no hiperplano
que \emph{melhor aproxima} $r$ em erro quadrático médio, sob a distribuição de $\X$ em
que os dados foram colhidos. É um resumo global, e útil como tal.

\emph{O que não se pode ler nele}, em três sentidos.
\begin{itemize}[leftmargin=1.2cm]
  \item \emph{Não é a inclinação de $r$.} Com $r$ curva, $\partial r/\partial x_1$ muda
        de ponto para ponto, e $2{,}3$ não é o valor dela em lugar nenhum em
        particular. Com várias covariáveis, pode nem ter relação com ela. Tome
        $X_2\sim N(0,1)$, $X_1=X_2^2+\varepsilon$, com $\varepsilon\sim N(0,s^2)$
        independente, e $r(\x)=x_2^2$: $r$ nem depende de $x_1$, e ainda assim,
        regredindo em $(1,x_1,x_2)$,
        \[
          \beta_1^\ast=\frac{\Cov(X_1,X_2^2)}{\Var(X_1)}=\frac{2}{2+s^2}\neq0 .
        \]
        A reta usa $x_1$ como \emph{procurador} da parte curva de $r$ que ela não
        consegue representar em $x_2$. (A conta: $X_2$ não é correlacionada nem com
        $X_2^2$ nem com $X_1$, pois $\E[X_2^3]=0$; então o coeficiente de $x_2$ é zero,
        o de $x_1$ é o da regressão simples, e $\Var(X_2^2)=\E[X_2^4]-1=2$.)
  \item \emph{Depende de onde os dados caíram.} Mesmo com uma covariável só: para
        $r(x)=x^2$ e $X\sim\mathrm{Unif}(0,a)$, a inclinação ótima é
        $\Cov(X,X^2)/\Var(X)=(a^3/12)/(a^2/12)=a$. Colhendo dados em $[0,1]$ o colega
        veria $1$; em $[0,2]$, veria $2$ --- com $r$ exatamente a mesma. Um número que
        muda com a faixa amostrada não descreve um mecanismo.
  \item \emph{"Aumentar $x_1$ aumenta $Y$" é uma frase causal}, sobre intervir no
        sistema. Nenhuma regressão a autoriza sozinha, nem com $r$ linear: comparar
        indivíduos que \emph{diferem} em $x_1$ não é o mesmo que \emph{mudar} o $x_1$
        de alguém.
\end{itemize}
\end{solucao}
